\section{Introduction}

Many scientific interventions change a law of motion rather than a scalar treatment. Diffusivity can be altered by a material, a boundary protocol can change a transport operator, and a spatial policy can modify the generator of a stochastic field. In these settings the object of interest is not merely a parameter. It is a post-intervention distribution of a functional of the field.

We study the linear parabolic model
\begin{equation}\label{eq:intro-spde}
 dU_t=\theta A_0U_t\,dt+\sigma Q^{1/2}dW_t,
 \qquad U_0=u_0,
\end{equation}
where $A_0$ is self-adjoint, negative, and has compact resolvent. The primary intervention is
\[
 \doop(\theta=\widetilde\theta),
\]
which replaces $\theta A_0$ by $\widetilde\theta A_0$ while leaving the other structural modules invariant. For a space-time weight $w$, the target is
\begin{equation}\label{eq:intro-target}
 F_{\widetilde\theta,w}(z;\vartheta)
 =\Pp_{\vartheta}\left\{
 \int_0^T\inner{w_t}{U_t^{\widetilde\theta}}\,dt\le z
 \right\},
 \qquad \vartheta=(\theta,\sigma).
\end{equation}

The paper addresses four distinct questions. First, what makes the formal replacement of an operator a structural intervention rather than a relabelled stochastic equation? Second, when is the counterfactual law identified from a factual observation experiment? Third, how can it be estimated uniformly over interventions, weights, and thresholds? Fourth, how do statistical and numerical errors interact at the inferential scale?

Our construction extends the finite-dimensional intervention logic of \citet{SokolHansen2014} to a diagonal parabolic SPDE. The post-intervention field is represented as the limit of intervened Galerkin structural causal models using the same Brownian path segments in every world. This common-noise convention is stronger than specifying marginal laws; it fixes the primitive disturbances that remain invariant under the intervention.

The statistical experiment uses exact discrete observations of the first $N$ spectral coordinates. This setting is close to the spectral inference literature surveyed by \citet{Cialenco2018}, but the target and asymptotic geometry are different. The exact transition likelihood produces two rates. With $\Lambda_N=\sum_{k\le N}\lambda_k$,
\[
 \sqrt{\Lambda_N}(\widehat\theta-\theta_0)
 \quad\text{and}\quad
 \sqrt{Nn}(\widehat\sigma-\sigma_0)
\]
have non-degenerate limits. Fixed absolute-intervention laws overwrite the factual drift and are therefore estimated at the faster scale rate. Factual and factual-versus-counterfactual effects retain the slower drift uncertainty.

The main contributions are as follows.

\begin{enumerate}[label=(\roman*)]
\item We prove exact finite-dimensional intervention semantics and uniform convergence of intervened Galerkin models to the post-intervention SPDE.
\item We give a fibre criterion for target identification and explicit finite-sample identification results under exact and noisy spectral observations.
\item We derive the exact likelihood, mixed Fisher-information limit, LAN expansion, efficient MLE, and moment convergence.
\item We establish uniform plug-in processes in $C(\cT)$, feasible multiplier bands, and contact-set tests for one-sided counterfactual hypotheses.
\item We prove operator-to-law derivatives for bounded perturbations and a separate analytic-semigroup theorem for the unbounded rescaling direction $A_0$.
\item We combine statistical, Galerkin, temporal, and operator approximation errors in one family-specific inference theorem.
\item We construct sharp non-identification examples. In particular, unrestricted feedback or open-loop forcing can absorb the factual operator and separate only after intervention.
\end{enumerate}

The contribution is deliberately narrower than causal discovery. The intervention is specified by the analyst. The theory identifies its target inside an admissible model class. It does not infer a causal graph from passive field observations, and it does not claim that every decomposition of an observed drift is structural.

\subsection*{Related literature}

Three literatures meet in this paper, and the delimitation matters.

\paragraph{Statistics for SPDEs.}
The spectral approach to drift estimation under continuous observation goes
back to \citet{HuebnerRozovskii1995}; see \citet{Lototsky2009} and the survey
of \citet{Cialenco2018}. An active recent strand studies discretely observed
parabolic equations: power variations and volatility estimation under
high-frequency sampling \citep{BibingerTrabs2020,Chong2020}, joint estimation
from space--time grids
\citep{HildebrandtTrabs2021,KainoUchida2021,TonakiKainoUchida2023,TonakiKainoUchida2024},
drift estimation from discretely sampled Fourier modes
\citep{CialencoDelgadoVencesKim2020}, local measurements
\citep{AltmeyerReiss2021}, and semilinear extensions
\citep{PasemannStannat2020}. Relative to this literature our observation
scheme is deliberately the cleanest one --- exact discrete observation of $N$
spectral coordinates --- but the output is of a different kind: a
\emph{joint} local asymptotic normality expansion with anisotropic diagonal
normalisation $\diag(\sqrt{\Lambda_N},\sqrt{Nn})$, efficiency and moment
convergence of the exact conditional MLE, pushed through to uniform
inference on post-intervention laws of field functionals. To our knowledge
both the mixed-rate LAN statement and any inferential theory for operator
interventions in SPDEs are new; the nearest discrete-observation results
concern estimation rates rather than experiment-level asymptotics.

\paragraph{Causal models for dynamical systems.}
Interventions in ODE systems were formalised as limits of structural
equation models by \citet{MooijJanzingScholkopf2013}, and extended to dynamic
equilibria and random differential equations by \citet{RubensteinEtAl2018}
and \citet{BongersMooij2018}; \citet{SokolHansen2014} treat SDEs via Euler
schemes, and \citet{PetersBauerPfister2022} survey the programme.
\Cref{thm:semantics} is the infinite-dimensional counterpart: the
post-intervention SPDE is the limit of intervened Galerkin structural causal
models driven by \emph{common} Brownian segments, so the intervention is
module surgery in the sense of \citet{Haavelmo1944} rather than a
relabelling of stochastic equations. The non-identification constructions of
\Cref{sec:nonid} calibrate exactly which invariances carry identifying
content.

\paragraph{Uniform inference on distributional targets.}
The plug-in processes and feasible bands use the functional delta method
\citep{vanDerVaart1998,vanDerVaartWellner1996}, multiplier approximations
for suprema \citep{ChernozhukovChetverikovKato2013}, directional
differentiability \citep{FangSantos2019}, and contact-set logic
\citep{AndrewsSoares2010,LintonSongWhang2010}. The finite-rank limit
geometry --- every target process is a deterministic derivative surface
loading a two-dimensional Gaussian --- makes these tools exact rather than
high-dimensional approximations, and the mixed rates propagate to
target-specific band widths. In econometric terms the targets are
counterfactual distributions and quantiles
\citep{ChernozhukovFernandezValMelly2013}; identification from discretely
sampled continuous-time dynamics connects to the aliasing literature
(\Cref{rem:no-aliasing}).

The paper is organised as follows. \Cref{sec:model} states the model and Gaussian target law. \Cref{sec:semantics} proves the intervention semantics. \Cref{sec:identification} treats identification. \Cref{sec:lan} develops the statistical experiment. \Cref{sec:uniform} gives uniform inference and bootstrap procedures. \Cref{sec:sensitivity} treats operator response and total error. \Cref{sec:nonid} gives the impossibility boundary. \Cref{sec:numerics} reports numerical validation. \ifdefined\MERGEDDOC Complete proofs of every statement are collected in \Cref{sec:proofs} (parts S0--S8), in dependency order and preceded by part S0, which restates every imported external theorem as an attributed local proposition whose hypotheses are verified at each point of use.\else Complete proofs of every statement are collected in the accompanying consolidated technical supplement (parts S0--S8), which supersedes the separate work-package archives; part S0 restates every imported external theorem as an attributed local proposition whose hypotheses are verified at each point of use.\fi

\subsection*{Notation}
$H=L^2(D)$ with inner product $\inner\cdot\cdot$ and norm $\norm\cdot$; $A_0$ is self-adjoint, negative, with compact resolvent, eigenpairs $(-\lambda_k,e_k)$, $0<\lambda_1\le\lambda_2\le\cdots$, and $\lambda_k\asymp k^{2/d}$ (Weyl growth). $\Lambda_N:=\sum_{k\le N}\lambda_k$. The parameter is $\vartheta=(\theta,\sigma)$, its true value $\vartheta_0$, and $A_{N,n}:=\diag(\sqrt{\Lambda_N},\sqrt{Nn})$ is the information normalisation. $\doop(\theta=\tilde a)$ denotes module surgery on the generator coefficient. Weak convergence is $\Rightarrow$; $\op(1),\Op(1)$ are stochastic order symbols; $C,c$ denote finite positive constants that may change between displays and whose dependencies are stated where they matter; $a\wedge b=\min(a,b)$, $a\vee b=\max(a,b)$, and $x_+=\max(x,0)$. For a compact index set $\cT$, $C(\cT)$ carries the supremum norm $\norm\cdot_\infty$. \ifdefined\MERGEDDOC\else \Cref{tab:crosswalk} maps each statement to its proof.\fi

\ifdefined\MERGEDDOC\else
\begin{table}[t]
\centering\small
\caption{Statement-to-proof crosswalk\ifdefined\MERGEDDOC\ (\Cref{sec:proofs})\else\ (technical supplement)\fi.}
\label{tab:crosswalk}
\begin{tabular}{ll}
\toprule
Main-text statement & Proof location \\
\midrule
\Cref{prop:gaussian-law} & Part S2 (Gaussian post-intervention law)\\
\Cref{prop:exact-scm,thm:semantics} & Part S1 (exact representation; Galerkin limits)\\
\Cref{thm:fibre,thm:exact-id,prop:partial} & Part S2 (identification)\\
\Cref{thm:lan} & Part S3 (LAN, efficiency, moment convergence)\\
\Cref{thm:processes,thm:quantile-tail} & Part S4 (uniform target processes)\\
\Cref{thm:bands} & Part S5 (multiplier, floor bands, contact sets)\\
\Cref{thm:sensitivity} & Part S6 (operator sensitivity)\\
\Cref{thm:total-error} & Part S7 (total error)\\
\Cref{thm:gauge-main,thm:nonid-assembly} & Part S8 (non-identification)\\
\bottomrule
\end{tabular}
\end{table}
\fi
